\begin{frame}{첫방문이 실제로 ``한번''이 아님을 보는 예}
\begin{center}\begin{varwidth}{0.92\textwidth}\usebeamercolor[fg]{normal text}
  
  \begin{tabular}{@{}clll@{}}
    \toprule
    스텝 & 위치 & 보상 $R$ & 리턴 $G$ \\
    \midrule
    0 & $(0,0)$ & $-0.1$ & $-0.139$ \\
    1 & $(0,1) \to$ (미끄러짐) $(0,0)$ & $-0.1$ & $-0.043$ \\
    2 & $(0,0)$ (두 번째 방문) & $-0.1$ & $+0.063$ \\
    3 & $(0,1)$ & $-0.1$ & $+0.181$ \\
    4 & $(0,0)$ (세 번째, 다시 미끄러짐) & $-0.1$ & $+0.312$ \\
    5$\sim$8 & (정상 루프 $\to$ 터미널) & $\cdots$ & $\to +1.0$ \\
    \bottomrule
  \end{tabular}
  \vskip0.8em
  \textbf{첫방문 MC는 첫 방문 리턴만 사용}
  \vskip0.3em
\end{varwidth}\end{center}
\note{\begin{itemize}
\item 같은 $(0,0)$ 에서 시작, 이번에는 미끄러짐이 몇 번:
  \small
\item $(0,0)$ 는 이 에피소드에서 \textbf{세 번} 방문, 매 방문마다
  리턴이 달랐다 ($-0.139,\ +0.063,\ +0.312$). 첫방문 MC 는
  \textbf{$-0.139$ 만} 한 번 더치고 나머지 두 방문은 무시한다 ---
  ``이 에피소드가 $(0,0)$ 에 대해 말하는 것''을
  \textbf{첫 방문 시점의 리턴 하나로 압축}.
\item (마지막 방문을 쓰면 $+0.312$, 세 방문 모두 쓰면
  $(-0.139+0.063+0.312)/3 = +0.079$ --- \textbf{같은 에피소드에서
  세 가지 다른 추정치}가 나온다.)
\end{itemize}}
\end{frame}
